Acre & 21,9 & 29\% & 9,7 a 11,2 & 1,9 a 2,2 & 36 \\
Alagoas & 19,2 & 73\% & 6,5 a 7,5 & 2,5 a 2,9 & 22 \\
Amapá & 21,2 & 31\% & 9,3 a 10,7 & 1,9 a 2,3 & 31 \\
Amazonas & 15,0 & 48\% & 8,2 a 9,4 & 1,5 a 1,8 & 37 \\
Bahia & 18,1 & 83\% & 10,8 a 12,5 & 1,4 a 1,7 & 51 \\
Ceará & 16,4 & 97\% & 11,0 a 12,6 & 1,2 a 1,5 & 30 \\
Distrito Federal & 19,2 & 23\% & 11,2 a 12,9 & 1,4 a 1,7 & 52 \\
Espírito Santo & 19,8 & 61\% & 11,6 a 13,3 & 1,4 a 1,7 & 27 \\
Goiás & 13,1 & 47\% & 8,6 a 9,9 & 1,2 a 1,5 & 20 \\
Maranhão & 19,3 & 93\% & 8,9 a 10,2 & 1,8 a 2,2 & 45 \\
Mato Grosso & 17,9 & 39\% & 9,6 a 11,0 & 1,6 a 1,9 & 43 \\
Mato Grosso do Sul & 18,6 & 49\% & 10,8 a 12,4 & 1,4 a 1,7 & 47 \\
Minas Gerais & 11,2 & 42\% & 12,6 a 14,4 & 0,7 a 0,9 & 58 \\
Pará & 18,2 & 71\% & 7,4 a 8,5 & 2,1 a 2,5 & 34 \\
Paraíba & 17,9 & 74\% & 6,4 a 7,4 & 2,3 a 2,8 & 31 \\
Paraná & 15,9 & 40\% & 9,9 a 11,4 & 1,3 a 1,6 & 37 \\
Pernambuco & 16,4 & 71\% & 8,2 a 9,4 & 1,7 a 2,0 & 23 \\
Piauí & 18,7 & 87\% & 10,3 a 11,9 & 1,5 a 1,8 & 48 \\
Rio Grande do Norte & 18,4 & 57\% & 8,3 a 9,5 & 1,9 a 2,2 & 21 \\
Rio Grande do Sul & 12,2 & 39\% & 11,5 a 13,2 & 0,9 a 1,1 & 30 \\
Rio de Janeiro & 16,1 & 73\% & 13,7 a 15,8 & 1,0 a 1,2 & 39 \\
Rondônia & 17,7 & 42\% & 6,9 a 7,9 & 2,2 a 2,6 & 25 \\
Roraima & 19,2 & 32\% & 6,6 a 7,5 & 2,5 a 2,9 & 29 \\
Santa Catarina & 13,6 & 44\% & 11,9 a 13,7 & 0,9 a 1,1 & 47 \\
São Paulo & 19,0 & 41\% & 13,0 a 15,0 & 1,2 a 1,5 & 45 \\
Sergipe & 16,4 & 55\% & 8,8 a 10,1 & 1,6 a 1,9 & 34 \\
Tocantins & 19,2 & 49\% & 11,3 a 13,0 & 1,4 a 1,7 & 38 \\
\midrule
\textbf{Média dos estados} & \textbf{17,4} & & \textbf{9,7 a 11,2} & \textbf{1,6 a 1,9} & \textbf{36} \\
